\documentclass[10pt,a4paper]{amsart}
\usepackage[margin=19mm]{geometry}
\usepackage{amsmath,amssymb,mathtools}
\usepackage[scr=rsfs,cal=euler]{mathalfa}
\usepackage{xfrac}
\usepackage[unicode,hidelinks,bookmarks=false]{hyperref}
\newcommand{\ie}{i.e.\ }
\newcommand{\cf}{cf.\ }
\newcommand{\resp}{respectively\ }
\newcommand{\Subsection}{\subsection}
\providecommand{\nobreakdash}{\nobreak\hbox{-}}
\setlength{\emergencystretch}{2em}
\multlinegap0pt
\usepackage{enumerate}
\usepackage{cancel}
\usepackage{amssymb}
\usepackage{xcolor}
\newtheorem{lem}{Lemma}
\newcommand{\E}{\mathbb{E}}
\newcommand{\Q}{\mathbb{Q}}
\newcommand{\R}{\mathbb{R}}
\newcommand{\Span}{\text{span}}
\newcommand{\WP}{\R\times_{\rho}\Q^n_{\epsilon}}
\newcommand{\al}{\alpha}
\newcommand{\co}{\colon}
\newcommand{\dt}{\partial_t}
\newcommand{\e}{\epsilon}
\newcommand{\nb}{\overline\nabla}
\newcommand{\np}{\nabla^\perp}
\newcommand{\nt}{\widetilde\nabla}
\newcommand{\pid}{\rangle}
\newcommand{\pie}{\langle}
\title{Surfaces with parallel mean curvature in warped product spaces}
\author{Fernando Manfio, Ver\^onica Reis, Feliciano Vit\'orio}
\date{Source: arXiv:2603.01394v1 (CC BY 4.0)}
\begin{document}
\maketitle

\noindent\emph{Source and licence.} Surfaces with parallel mean curvature in warped product spaces. Version arXiv:2603.01394v1 (CC BY 4.0), distributed by arXiv under the Creative Commons Attribution 4.0 International licence, http://creativecommons.org/licenses/by/4.0/, as recorded in the arXiv licence metadata for this exact version and shown on the abstract page https://arxiv.org/abs/2603.01394v1. Full licence notice: \texttt{licenses/arxiv-2603.01394-CC-BY-4.0.txt}.\par\smallskip

\noindent\textit{Editorial scope.} Counted unit: the article's lem lem:RedCodim (source0.tex lines 263--265). The article's complete original proof of it is delivered with it (source0.tex lines 266--294, one \texttt{proof} environment of 29 lines). No mathematics is written, repaired or re-derived: the statement and the proof are the article's own lines, in the article's own notation, and no retained line is substituted or re-flowed.  Retained uncounted as the source-local context the proof needs: lines 139--142; lines 147--166; lines 237--252.  Nothing was omitted from any retained span, so the package carries no omission record.  No retained line contains a citation, so the wrapper carries no bibliography and no bibliography entry.  No retained line contains a figure environment, an image-inclusion command or drawing code, so no image is delivered and the package declares no asset dependency.  Editorial formatting only, in addition: an AMS \texttt{amsart} preamble that restates the article's own declarations and macros used by the retained lines, namely the environments \texttt{lem} on separate counters, together with the standard packages those lines need.  The retained environments print in this document as Lemma 1 in order of appearance, whereas the article prints the same items with the numbering of its own full document.  The complete native source of the article as distributed in the arXiv e-print for this exact version is bundled unchanged as \texttt{sources/195559/source0.tex}.
\begin{equation} \label{eq:Nabladelt}
\nb_Z\dt = \frac{\rho'(t)}{\rho(t)}\left(Z-\pie Z,
\dt\pid\dt\right),
\end{equation}

Given an isometric immersion $f\co M^m\to\WP$, a tangent vector field $T$ on $M^m$ and a normal vector field $\eta$ along $f$ are defined by
\begin{equation} \label{eq:eta}
\partial_t= f_{\ast} T+\eta.    
\end{equation}
Using \eqref{eq:Nabladelt}, we obtain by differentiating \eqref{eq:eta} that
\begin{eqnarray} \label{eq:derT}
\nabla_XT - A^f_\eta X = \frac{\rho'(t)}{\rho(t)}
\left(X-\langle X,T\rangle T\right)
\end{eqnarray}
and
\begin{eqnarray} \label{eq:dereta}
\al_f(X,T) + \nabla^{\perp}_X\eta = -\frac{\rho'(t)}{\rho(t)}
\langle X,T\rangle\eta,
\end{eqnarray}
for all $X\in TM$, where $\alpha_f$ and $\np$ denote the second fundamental form and the normal connection of $f$, respectively. Here $A^f_\eta$ stands for the shape operator of $f$ in the
direction $\eta$, given by
\[
\pie A^f_\eta X,Y\pid = \pie\al_f(X,Y,\eta\pid,
\]
for all $X,Y\in TM$.

\begin{eqnarray} \label{eq:shapes}
\begin{aligned}
\nt_Xi_\ast\xi &= i_\ast\nb_X\xi + \alpha_i(f_\ast X,\xi) \\
& = -\tilde{f}_\ast A^f_{\xi}X + i_\ast\np_X\xi+ \alpha_i(f_\ast X,\xi)
%\e\left(\frac{h'(t)}{\rho(t)}+\e\frac{\rho''(t)}{h'(t)}\right)\langle X, T\rangle \langle\eta,\xi\rangle N_t.
\end{aligned}
\end{eqnarray}
Therefore, from \eqref{eq:shapes} and using Weingarten formula, one has
\begin{eqnarray} \label{eq:A^fA^tf}
A^{\tilde f}_{i_{\ast}\xi}=A^f_{\xi}
\end{eqnarray}
and
\begin{eqnarray} \label{eq:nabtilperp}
\widetilde{\nabla}^{\perp}_{X} i_{\ast} \xi=  i_\ast\np_X\xi+ \e\left(\frac{h'(t)}{\rho(t)}+\e\frac{\rho''(t)}{h'(t)}\right)\langle X, T\rangle \langle\eta,\xi\rangle N_t,
\end{eqnarray}
for every $\xi\in TM^\perp_f$, where $\widetilde{\nabla}^{\perp}$ is the normal connection of $\tilde{f}$. 

\begin{lem} \label{lem:RedCodim}
Let $f\co M^m\to\WP$ be an isometric immersion and assume that $L:= N_{1}+\Span\{\eta\}$ is a subbundle of $TM^\perp_f$ of rank $l<n+1-m$ and that $\nabla^{\perp}N_1\subset L$. Then f reduces codimension to $l$.
\end{lem}

\begin{proof}
It follows from \eqref{eq:dereta} that $\np_X\eta\in L$, for every $X\in TM$. Since we have $\np N_1\subset L$ by assumption, it follows that $L$ is a parallel subbundle of $TM^\perp_f$.
%Consider now the composition $\tilde{f}=i\circ f$, where $i\co\WP\to\E^{n+2}$ is the inclusion, and 
Let $L^\perp$ denote the orthogonal complement of $L$ in $TM^\perp_f$. Given $\xi\in L^\perp=N_1^\perp\cap\{\eta\}^\perp$, from \eqref{eq:nabtilperp} and the fact that $L$ is a parallel subbundle of $TM^\perp_f$, we obtain
\[
\widetilde{\nabla}_X^\perp i_{\ast}\xi = 
i_{\ast}\nabla^{\perp}_{X} \xi \in i_{\ast}L^{\perp},
\]
where $\widetilde{\nabla}$ denotes the derivative of $\E^{n+2}$. This shows that $i_{\ast}L^{\perp}$ is a parallel subbundle of $TM^\perp_{\tilde f}$. On the other hand, it follows from \eqref{eq:A^fA^tf} that
\[
i_{\ast}L^{\perp}\subset\widetilde{N}_{1}^{\perp},
\]
where $\widetilde{N}_1(x)$ is the first normal space of $\tilde{f}$ at $x\in M^m$. We obtain from the Weingarten formula for $\tilde{f}$ that
\[
\widetilde{\nabla}_{X}i_{\ast}\xi = -\tilde{f}_{\ast}A_{ i_{\ast}\xi}^{\tilde{f}}X + \widetilde{\nabla}_{X}^{\perp}  i_{\ast}\xi=\widetilde{\nabla}_{X}^{\perp}  i_{\ast}\xi \in i_{\ast}L^{\perp},
\]
This implies that $i_{\ast}L^{\perp}$ is a constant subspace of $\E^{n+2}$, which is orthogonal to $\dt$. 
%In fact, given $\xi \in i_{\ast}L^{\perp}$, we have $\langle \xi, f_{\ast} T+\eta\rangle= \langle\xi,\eta\rangle= 0$.
Denote by $K$ the orthogonal complement of $i_{\ast}L^{\perp}$ in $\E^{n+2}$. For any fixed point $x_0\in M^m$, we have
\[
\tilde f(M)\subset\tilde{f}(x_0)+K.
\]
Since $K$ contains $\dt$ and $N_t(x_0)$, it also contains the position vector $\tilde{f}(x_0)$. Thus, we have $\tilde{f}(x_0)+K=K$. We conclude that
\[
\tilde{f}(M)\subset(\R\times_{\rho}\Q^{n}_{\epsilon})\cap K
= \R\times_{\rho}\Q^{m+l-1}_{\epsilon},
\]
and this completes the proof.
\end{proof}

\end{document}
